\documentclass[11pt,a4paper]{article}
\usepackage[margin=1.8cm]{geometry}
\usepackage{amsmath,amssymb,array,tabularx,fontspec,cancel}
\setmainfont{FreeSerif}
\setlength{\parindent}{0pt}\setlength{\parskip}{0pt}
\renewcommand{\arraystretch}{1.15}
\newcounter{sq}
\newcommand{\sq}[1]{\stepcounter{sq}\par\smallskip\noindent\hangindent=11em\hangafter=1
  \makebox[11em][l]{\textbf{\Roman{sq}.}}#1\par}
\newcommand{\dtab}[2]{\par\smallskip\noindent\hspace*{11em}\begin{tabular}{#1}\hline #2\\ \hline\end{tabular}\par\smallskip\hspace*{11em}}
\begin{document}
\begin{center}
  {\Huge\bfseries The Hope Science Academy}\\[3pt]
  {\small Excellence is our standard not goal}\\[2pt]
  {\small 247-E-1, (Opposite LDA School) Johar Town, Lahore. PH: 0300-8194789, 0322-4157001}
\end{center}
\medskip\noindent
\begin{tabularx}{\linewidth}{|X|X|X|X|}
\hline
\textbf{Student Name:} & \textbf{Roll Num:} & \textbf{Class:} 12th & \textbf{Paper Code:} \\ \hline
\textbf{Subject Name:} Business Statistics & \textbf{Time Allowed:} & \textbf{Total Marks:} 658 & \textbf{Exam Date:} 02-Nov-2026 \\ \hline
\multicolumn{4}{|l|}{\textbf{Exam Syllabus:} FULL BOOK} \\ \hline
\end{tabularx}
\bigskip
\noindent\textbf{Q2. Write short answers of the following questions.}\hfill\textbf{1$\times$311=311}

\sq{Give Italian, Latin and German words used for the term statistics.}
\sq{What is statistics?}
\sq{Define statistics in your own words.}
\sq{Explain the term statistics in two different meanings.}
\sq{What are the two branches of statistics?}
\sq{List the phases of statistical enquiry.}
\sq{Differentiate census from sample survey.}
\sq{Why sample survey is preferred over census?}
\sq{Distinguish a complete enumeration from sample survey.}
\sq{Describe two limitations of statistics.}
\sq{Describe two limitations of statistics.}
\sq{What is variable?}
\sq{What is constant?}
\sq{Define discrete variable and continuous variable.}
\sq{Define discrete variable by giving examples.}
\sq{What is quantitative variable?}
\sq{What is qualitative variable?}
\sq{What is population?}
\sq{What is sample?}
\sq{Define population and sample.}
\sq{What is a parameter?}
\sq{What is a statistic?}
\sq{What is descriptive statistics?}
\sq{What is inferential statistics?}
\sq{What is data?}
\sq{What is primary data?}
\sq{What is secondary data?}
\sq{Give two sources of secondary data.}
\sq{Define continuous data and give two examples.}
\sq{Define attribute and by giving examples.}
\sq{Note major sources for collecting secondary data.}
\sq{Name five fields of application of statistics.}
\sq{Differentiate primary data from secondary data.}
\sq{Name the methods of collecting primary data.}
\sq{Describe editing of data.}
\sq{Why editing of data is necessary?}
\sq{Define constant with an example.}
\sq{What is questionnaire?}
\sq{Write down the applications of statistics in business.}
\sq{Give the two characteristics of statistics.}
\sq{Give an example each of primary and secondary data.}
\sq{What are the two sources of statistical data?}
\sq{Differentiate between descriptive and inferential statistics.}
\sq{Write down three limitations of Statistics.}
\sq{Define continuous variable with example.}
\sq{Define primary data.}
\sq{Define ``Statistics'' in singular sense.}
\sq{Briefly describe the important of statistics in administration.}
\sq{What is continuous variable?}
\sq{What is qualitative variable?}
\sq{What are the various sources of primary data collection?}
\sq{Define discrete variable.}
\sq{What is a constant and variable?}
\sq{What is a parameter?}
\sq{Define primary data.}
\sq{What are the main sources of collecting secondary data?}
\sq{Describe the method of direct personal investigation for collecting data.}
\sq{Describe the method of registration through local sources for collecting data.}
\sq{What is classification?}
\sq{Define tabulation.}
\sq{What are the different parts of a statistical table?}
\sq{Draw sketch of a table indicating its necessary parts.}
\sq{What is ``frequency distribution''?}
\sq{What is an array?}
\sq{What is meant by grouped data?}
\sq{Define class limits?}
\sq{Define class interval.}
\sq{Define class mark.}
\sq{What steps should we consider to make a frequency distribution?}
\sq{What is a relative frequency distribution?}
\sq{What is meant by graph?}
\sq{Define Histogram.}
\sq{What is frequency polygon?}
\sq{What are the different types of charts used in presenting a statistical data?}
\sq{When do we use simple bar chart, multiple bar chart, component bar chart and percentage component bar chart?}
\sq{What is it the importance of using multiple bar chart?}
\sq{When do you prefer to use component bar chart?}
\sq{When should you prefer to use pie chart?}
\sq{When is it to use percentage component bar chart?}
\sq{When do you prefer to draw historigram?}
\sq{What is a pie chart?}
\sq{What are the important factors in drawing a graph?}
\sq{Give two advantages of graphs.}
\sq{Name two parts of statistical table.}
\sq{What is Historigram?}
\sq{Name the types of graphs.}
\sq{Define class interval.}
\sq{What do you understand by the term ``Classification''?}
\sq{Define multiple bar diagram.}
\sq{Define frequency polygon.}
\sq{What is meant by class boundaries?}
\sq{What are main classes of graphs?}
\sq{Define pie-chart.}
\sq{Enlist main parts of table.}
\sq{Distinguish between histogram and historigram.}
\sq{Differentiate between ungrouped and grouped data.}
\sq{Differentiate between ungrouped and grouped data.}
\sq{Differentiate between class limits and class boundaries.}
\sq{List main parts of a table.}
\sq{What is frequency distribution?}
\sq{What is difference between diagrams and graphs?}
\sq{What is meant by one-way classification?}
\sq{Define the term class mark.}
\sq{Define the term ungroup data.}
\sq{Define class interval.}
\sq{Describe the four bases of classification of data.}
\sq{Define tabulation.}
\sq{What is pie chart?}
\sq{What is classification?}
\sq{What is an array?}
\sq{Define class mark.}
\sq{Define histogram.}
\sq{What are the different types of charts used in presenting of statistical data?}
\sq{Differentiate between group and ungroup data.}
\sq{Define frequency polygon.}
\sq{Define class interval.}
\sq{What is ``Measure of Central Tendency''?}
\sq{What is an average?.}
\sq{Describe the qualities of a good average.}
\sq{Name different averages and define one of them.}
\sq{Define Arithmetic Mean.}
\sq{Explain the ``Direct Method'' for finding the arithmetic mean.}
\sq{What do you understand by ``Change of Origin''?}
\sq{Give two properties of ``Arithmetic Mean''.}
\sq{Explain the concept of ``Weighted Arithmetic mean''}
\sq{Give merits and demerits of ``Arithmetic Mean''.}
\sq{Define Median.}
\sq{What purpose do you think an average serve?}
\sq{Compare strengths and weaknesses of mode.}
\sq{Compare advantages and disadvantages of median.}
\sq{List merits of arithmetic mean.}
\sq{Give two advantages of A.M.}
\sq{Define mode with examples.}
\sq{Write two merits of median.}
\sq{Write two properties of arithmetic mean.}
\sq{How do you differentiate ordinary mean from weighted mean?}
\sq{Write empirical relationship among mean, median and mode.}
\sq{Draw a frequency curve showing positively skewed distribution property indicating positions of mean, median and mode.}
\sq{Explain relationship among mean, median and mode by i) equation and ii) sketch.}
\sq{For a certain frequency distribution, the value of mean is 15 and median is 20. What will be the value of mode?}
\sq{Write desirable qualities of a good average.}
\sq{For a moderately skewed distribution mode=60 and mean 30. Find median.}
\sq{The mean is affected by charge of origin and scale, explain.}
\sq{Write two common advantages and disadvantages of median and mode.}
\sq{A curve whose tail is longer to the right, what you would call it?}
\sq{Find out the missing value? i) mode=3 Median-? Mean ii) Median=? (2 mean + mode)}
\sq{If for 10 observations, $\sum(x-23)=-17$, then find the value of mean.}
\sq{The following data shows number of absent students during the month of November 2011 from the I.Com Classes. Make a frequency distribution table class interval as one. 3,4,5,6,7,1,0,2,3,4,5,7,8,7,2,1,5,6,7,8,9,10,6,7,3}
\sq{The following data indicate the size of shoes sold at a store during the given weeks. Find model size of shoes.\dtab{|l@{\hspace{1.2em}}c@{\hspace{1.2em}}c@{\hspace{1.2em}}c@{\hspace{1.2em}}c@{\hspace{1.2em}}c@{\hspace{1.2em}}c@{\hspace{1.2em}}c@{\hspace{1.2em}}c|}{Size of Shoes & $4\frac{1}{2}$ & 5 & $5\frac{1}{2}$ & 6 & $6\frac{1}{2}$ & 7 & $7\frac{1}{2}$ & 8 \\ No. Pairs Sold & 2 & 5 & 7 & 20 & 25 & 17 & 15 & 4}}
\sq{Write desirable qualities of a good average.}
\sq{In a certain distribution mean is 45 and median is 50 then find mode. Also give the assumption you take with reference the shape of the distribution.}
\sq{What is meant by measure of central tendency?}
\sq{Given $D=X-2025$, $\sum fD=-10750$ and $\sum f=500$, find arithmetic Mean.}
\sq{Define Arithmetic Mean.}
\sq{If $l=28$, $f_{m}=25$, $f_{1}=20$, $f_{2}=18$ and $h=7$, compute mode}
\sq{Write three merits of Arithmetic Mean.}
\sq{The minimum temperature in Murree for the first 10 days of March was -1,-2,0,3,3,4,3,2,6,1. Find Median.}
\sq{Define weighted arithmetic mean.}
\sq{The mean of 100 items is 5 and that of 150 items is 4.8. Find combined mean.}
\sq{In moderately skewed distribution mean=25 and mode=28. Find the value of median.}
\sq{When weighted average is required?}
\sq{Why mode is a positional value?}
\sq{Write the quality of a good average.}
\sq{Define mode with examples.}
\sq{Give the advantages of arithmetic mean.}
\sq{Define median.}
\sq{if $=19.5$, $h=5$, $f_{m}=25$, $f_{1}=15$ and $f_{2}=20$, then find mode?}
\sq{If $=19.5$, $h=5$, $f_{m}=25$, $f1=15$ and $f2=20$, then find mode?}
\sq{Give the formula of coding method to compute arithmetic mean.}
\sq{Define arithmetic mean.}
\sq{Write two demerits of mode.}
\sq{Define average.}
\sq{Find the median of the data 5, 6, 9,5,4,3,8,7,6.}
\sq{For a skewed distribution, median =40 and mean=30, find mode.}
\sq{If $\sum x=450$ and $n=20$, find arithmetic mean?}
\sq{Find A.M given that $x=10+5u$, $\sum fu=46$ and $n=125$.}
\sq{Find out Arithmetic Mean given $\sum x=308$ and $n=7$.}
\sq{What is an index number?}
\sq{What are the types of an index number?}
\sq{What is the base period?}
\sq{What is price index number?}
\sq{What is quantity index number?}
\sq{What is simple index number?}
\sq{List the points you will take into consideration for construction of index number.}
\sq{Define composite index number.}
\sq{Define simple and composite index numbers.}
\sq{Define chain base method.}
\sq{Define weighted index number.}
\sq{What is un-weighted index number?}
\sq{What is weighted index number?}
\sq{What are the methods available for computing simple index number?}
\sq{Define price relative.}
\sq{Define link relative.}
\sq{Define chain indices.}
\sq{What are the formulas of the weighted price index numbers?}
\sq{What do you understand by base period \& how it is selected?}
\sq{Define Laspeyre’s index number.}
\sq{Define Paasche’s index number.}
\sq{What are the main steps involved in the construction of an index number?}
\sq{Note down four uses of the index numbers.}
\sq{What are the limitations of index numbers?}
\sq{Laspeyre’s index =115, Paasche’s index= 111, Fisher’s index =?}
\sq{Laspeyre’s index = 115, Fisher’s index = 112.98, Paasche’s index=?}
\sq{What will be Laspeyre’s index if Paasche’s index and Fisher’s index are 108 and 112 respectively?}
\sq{If product of Laspeyre’s and Paasche’s indices is 15376. Find Fisher’s index.}
\sq{Convert money price Rs 1000 for 40 Kg into quantity price.}
\sq{The quantity price 0.04 Kg per rupee is equivalent to what rupee per 40 Kg}
\sq{Write the formula given by Fisher.}
\sq{What is ``Fixed base method'' in construction of price index number?}
\sq{Differentiate between fix bas and chain base method.}
\sq{Find price relatives for the data given below, using average of last three years as base:\dtab{l@{\hspace{1em}}c@{\hspace{1em}}c@{\hspace{1em}}c@{\hspace{1em}}c@{\hspace{1em}}c@{\hspace{1em}}c@{\hspace{1em}}c@{\hspace{1em}}c@{\hspace{1em}}c}{Year & 1990 & 1991 & 1992 & 1993 & 1994 & 1995 & 1996 & 1997 & 1998 \\ Price & 15 & 16 & 20 & 22 & 24 & 23 & 25 & 15 & 20}}
\sq{Define an index number.}
\sq{Define chain indices.}
\sq{What is the base period?}
\sq{Define link relative.}
\sq{What are the uses of index numbers?}
\sq{What is the formula of value index number?}
\sq{Given $\sum p_{n}q_{n}=1356,\ \sum p_{o}q_{o}=1600$ then find current year weighted index.}
\sq{Define simple index number.}
\sq{Define link relative.}
\sq{Explain fixed base method.}
\sq{Write any two uses of index number.}
\sq{Define price index number}
\sq{Define quantity index number.}
\sq{Define chain base method.}
\sq{Enlist any two uses of index numbers}
\sq{If $\sum p_{n}q_{n}=2560$ and $\sum p_{o}q_{o}=2173$ then find the value of Paasche's price index number?}
\sq{Given $\sum IW=14825$ and $\sum W=115$ then find weighted average of relative's price index number.}
\sq{Define price relatives.}
\sq{Define the term weighted index.}
\sq{Give W=8500 and [formula missing in the source], construct consumer price index number.}
\sq{Define the link relative.}
\sq{Given N=5 and $\sum\left(\dfrac{P_{n}}{P_{o}}\right)\times100=510$ find price index by average of price relative method.}
\sq{If $\sum p_{1}q_{1}=481$ and $\sum p_{o}q_{1}=402$ then find current year weighted index number?}
\sq{Define a set}
\sq{What are the elements of a set?}
\sq{How will you define finite and infinite set?}
\sq{What is a subset? What are its types?}
\sq{What do you understand by power set?}
\sq{Define the null set.}
\sq{What is meant by universal set?}
\sq{Define ``Venn diagram''.}
\sq{Define union of two sets A and B.}
\sq{Define intersection of two sets A and B.}
\sq{Define complement of a set.}
\sq{What do you understand by product set?}
\sq{What is ``n factorial''?}
\sq{What is the difference between permutation and combination?}
\sq{What is the relationship between permutation and combination?}
\sq{How would you find the difference between two sets?}
\sq{Write down the ``De-Morgan’s Law'' of sets.}
\sq{What do you understand by permutation with repetition?}
\sq{Define equivalent sets.}
\sq{What are identity laws of sets?}
\sq{What are important laws of sets?}
\sq{What are associative laws of sets?}
\sq{Write down commutative laws of sets.}
\sq{Write down distribution laws of sets.}
\sq{Write down complement laws of sets.}
\sq{What is meant by universal set?}
\sq{Define equivalent set.}
\sq{What is an empty set?}
\sq{Define an experiment, what are its types?}
\sq{Define a scientific experiment.}
\sq{Define a random experiment.}
\sq{Explain the characteristics of a random experiment.}
\sq{What is a trial?}
\sq{What is an outcome?}
\sq{Define sample space.}
\sq{Define sample point.}
\sq{Define an event, what are its types?}
\sq{Define simple event.}
\sq{Define compound event.}
\sq{Define the null (impossible) event.}
\sq{Define the sure event.}
\sq{What are mutually exclusive events (MEE)?}
\sq{What are not mutually exclusive events?}
\sq{What do you understand by collectively exhaustive events?}
\sq{Define equally likely events.}
\sq{What do you know about dependent events?}
\sq{Define independent events.}
\sq{Explain mutually exclusive and equally likely events.}
\sq{Write sample space for toss of a fair cubical dice.}
\sq{Write sample space for the experiment ``toss a pair of coins''.}
\sq{Give possible number of sample points if two items are selected from 10 distinct items.}
\sq{Find number of possible outcomes if a coin is thrown with a pair of cubical dice.}
\sq{Write two equally likely events in toss of a coin.}
\sq{What do you understand by random experiment?}
\sq{What is meant by not mutually exclusive events?}
\sq{Define conditional probability.}
\sq{Define mutually exclusive events.}
\sq{What is ``probability''?}
\sq{Make a ``Sample space'', when a cubical die is rolled.}
\sq{Two coins are tossed once. Find: i) Probability that two heads appeared. ii) Probability that exact one head appeared.}
\sq{A fair die is rolled once what is the probability of obtaining: i) Six ii) An even number iii) The number greater than 4}
\sq{List the sample space, when coin is tossed 3 times.}
\sq{What is the classical definition of probability?}
\sq{Distinguish between simple and compound events.}
\sq{In how many ways can six different books be arranged on?}
\sq{If A=B are two independent events and P(A)=0.40, P(B)=0.35, find $(A\cup B)=?$}
\sq{For any two events A and B $P(A)=\dfrac{2}{3}$, $P(A\cup B)=\dfrac{7}{12}$ and $P(A\cap B)=\dfrac{5}{12}$ find P(B)=?}
\sq{If one card is drawn from 52 playing cards, and the probability that it is a King of Heart Cards?}
\sq{Define a random experiment.}
\sq{Define sure events.}
\sq{Explain the term equally likely events.}
\sq{Define the term compound event.}
\sq{Differentiate between simple and compound events.}
\sq{Write sample space for toss of a fair cubical die.}
\sq{What is the probability of getting trail on a throw of coin?}
\sq{What is the probability of getting sum 7, when two dice thrown?}
\sq{State the addition law for two not mutually exclusive events}

\end{document}
